\section{Inherited terminal interface}
\label{sec:tpc42-interface}

We retain only the terminal data proved in TPC--36, TPC--39, and
TPC--41 \citep{WangTPC36,WangTPC39,WangTPC41}.  Fix
\(h\in\Z\setminus\{0\}\).  A physical source row is
\begin{equation}
 \alpha=(\ell_\alpha,d_\alpha),
 \qquad m_\alpha=\ell_\alpha d_\alpha\asymp Q,
 \qquad \ell_\alpha\asymp L_0,
 \qquad d_\alpha\asymp D_0,
 \qquad L_0D_0=Q,
 \qquad QJ\asymp X.
 \label{eq:tpc42-row-scales}
\end{equation}
The source separation makes \(m_\alpha\) determine \(\alpha\), and the
two inherited source systems obey
\begin{equation}
 \sum_\alpha |\gamma_\alpha^{(i)}|^2
 \ll X^{o(1)}Q,
 \qquad i=1,2.
 \label{eq:tpc42-source-l2}
\end{equation}
The orbit variable ranges through a set \(\cJ\) of \(O(J)\) positive
integers.

Choose three pairwise distinct auxiliary primes from the three disjoint
TPC--41 intervals, satisfying
\begin{equation}
 \max\cJ<q_i,
 \qquad
 \max_\alpha d_\alpha<q_i<\min_\alpha\ell_\alpha,
 \qquad q_i\nmid h,
 \qquad i=1,2,3,
 \label{eq:tpc42-admissible-primes}
\end{equation}
and put \(M=q_1q_2q_3\).  For sufficiently large \(X\), the size and
endpoint conditions give
\begin{equation}
 (j,M)=1\quad(j\in\cJ),
 \qquad
 \max_{\alpha,\beta}|m_\alpha-m_\beta|<M.
 \label{eq:tpc42-no-wrap-data}
\end{equation}
At the optimized endpoint,
\begin{equation}
 Q=X^{267/400+o(1)},
 \qquad J=X^{133/400+o(1)},
 \qquad M=X^{399/400+o(1)}.
 \label{eq:tpc42-endpoint}
\end{equation}

Let \(B\asymp X\) be the completed determinant window.  To avoid
confusing the source scale \(L_0\) with the alias half--width, write
\begin{equation}
 L_B=\left\lfloor\frac BM\right\rfloor,
 \qquad q=L_B+1,
 \qquad \cI_B=\{-L_B,\ldots,L_B\},
 \qquad K_B=2L_B+1.
 \label{eq:tpc42-alias-data}
\end{equation}
Then
\begin{equation}
 K_B=X^{1/400+o(1)}
 =X^{o(1)}\frac{Q}{J^2},
 \qquad
 \frac{Q}{K_B}=X^{o(1)}J^2.
 \label{eq:tpc42-missing-ledger}
\end{equation}

For \(k\in\cI_B\), define
\begin{align}
 n_{\alpha,j,k}
 &=m_\alpha j+h+kM,
 \label{eq:tpc42-terminal-n}\\
 b_{\alpha,j,k}
 &=-\mu(n_{\alpha,j,k})\log(n_{\alpha,j,k})
   \Phi_0(n_{\alpha,j,k}),
 \label{eq:tpc42-terminal-b}
\end{align}
with the second quantity interpreted as zero outside the positive
terminal support.  The left coefficient and fixed--alias output are
\begin{align}
 c^L_{\alpha,\gamma,j}
 &=\gamma_\alpha^{(1)}
   \mathfrak A_{\alpha,\gamma}^{\rm act}(j),
 \label{eq:tpc42-left-c}\\
 [\mathscr T_k^L]_{\gamma,j}
 &=\sum_\alpha c^L_{\alpha,\gamma,j}b_{\alpha,j,k}.
 \label{eq:tpc42-left-column}
\end{align}
The right output interchanges the row systems.  The multiplier contains
the actual row separation, gcd, residue, and smooth masks, and
\begin{equation}
 \sup_{\alpha,\gamma,j}
 |\mathfrak A_{\alpha,\gamma}^{\rm act}(j)|
 \ll X^{o(1)}.
 \label{eq:tpc42-multiplier-bound}
\end{equation}
In particular, \(c^L\) is independent of \(k\), while its source factor
contains \(\mu(d_\alpha)\).

Source separation and the dyadic row window also give the crude counts
\begin{equation}
 \#\{\alpha\}+\#\{\gamma\}\ll X^{o(1)}Q.
 \label{eq:tpc42-row-counts}
\end{equation}
No coprimality of \(m_\alpha j+h\) with \(M\) is assumed in the physical
identities below.  The scalar theorem of TPC--41 sees only the unit
residue sector; keeping all residue classes avoids an unrecorded
modulus-dependent pruning of the physical packet.

For the minimal \(q=L_B+1\) Fourier bank, put
\begin{align}
 u_{\alpha,j,0}&=b_{\alpha,j,0},
 \label{eq:tpc42-folded-u0}\\
 u_{\alpha,j,a}
 &=b_{\alpha,j,a}+b_{\alpha,j,a-q},
 \qquad 1\le a\le L_B.
 \label{eq:tpc42-folded-ua}
\end{align}
The left folded energy and its equal--row part are
\begin{align}
 \cE_{\rm fold,L}
 &=\sum_{\gamma,j}\sum_{a=0}^{L_B}
   \left|\sum_\alpha
   c^L_{\alpha,\gamma,j}u_{\alpha,j,a}\right|^2,
 \label{eq:tpc42-folded-energy}\\
 \cE_{\rm same,L}
 &=\sum_{\gamma,j}\sum_{a=0}^{L_B}\sum_\alpha
   |c^L_{\alpha,\gamma,j}|^2|u_{\alpha,j,a}|^2.
 \label{eq:tpc42-same-energy}
\end{align}
Writing \(\cF_L\) for the ordered \(\alpha\ne\beta\) cross--row sum,
\begin{equation}
 \cE_{\rm fold,L}
 =\cE_{\rm same,L}+\operatorname{Re}\cF_L.
 \label{eq:tpc42-folded-split}
\end{equation}
TPC--41 proves, together with the symmetric right estimate,
\begin{equation}
 \boxed{
 \cE_{\rm same,L}+\cE_{\rm same,R}
 \ll X^{o(1)}Q^2JK_B
 =X^{o(1)}\frac{Q^3}{J}.}
 \label{eq:tpc42-inherited-same-bound}
\end{equation}
Thus the remaining minimal--bank theorem is exactly the one--sided bound
\begin{equation}
 \operatorname{Re}(\cF_L+\cF_R)
 \ll X^{o(1)}Q^2JK_B.
 \label{eq:tpc42-inherited-gate}
\end{equation}
